\paragraph{Normal Forms}
A \emph{normal form} is typically defined with respect to a step relation, stating that each type cannot take a \emph{single} step:

\begin{code}[hide]
  module _ {Δ : KEnv} {κ : Kind} where 
    _—→ₜ_ : (_ _ : Type Δ κ) → Set
    _ —→ₜ _  = ⊤ 
\end{code}
\begin{code}
    _↛ₜ : (τ : Type Δ κ) → Set 
    τ ↛ₜ = ¬ (Σ[ υ ∈ (Type Δ κ) ] (τ —→ₜ υ))
\end{code}

\Ni However, we haven't such a relation, viewing any derivation of $\tau \TEqv \upsilon$ instead as a \emph{sequence} of steps. In turn, we instead define normal forms as fixed points of normalization:

\begin{code}
  NormalForm  : Type Δ κ → Set
  NormalForm τ = (⇑ ∘ ⇓) τ ≡ τ
\end{code}

\Ni In other words, normalization has no effect on a type in normal form. Stability implies this property for all normal types.

\begin{code}
  allNormal : (τ : NormalType Δ κ) → NormalForm (⇑ τ)
  allNormal τ = cong ⇑ (stability τ)
\end{code}